La estimación del rendimiento agrícola a escala de parcela plantea un problema de predicción en el que confluyen información ambiental heterogénea, observaciones espaciotemporales irregulares y procesos agronómicos parcialmente conocidos. La producción de cebada responde a la disponibilidad de agua, su distribución durante el ciclo, el desarrollo fenológico, las propiedades edáficas, el relieve y las prácticas de manejo. Tales efectos no necesariamente son constantes entre localidades o campañas. La literatura ha documentado la relación entre índices espectrales y rendimiento y la sensibilidad del cultivo al estrés hídrico en etapas específicas \parencite{Mirosavljevic2018,Sharifi2020,Chanev2025,Hajkova2019,Bello2022,Mogensen1980,Zhang2015,King2024}.

El Reto AgroCebada FIRA 2026 presenta un escenario delimitado: existen 197 parcelas con covariables observables, pero solamente 138 cuentan con un rendimiento registrado y las 59 restantes requieren estimación. La unidad estadística es la parcela identificada por \texttt{ID\_POLIGONO}, no cada fila satelital. El objetivo es reconstruir \texttt{RENDIMIENTO\_T\_HA}, medido en toneladas por hectárea, para un conjunto fijo de unidades correspondientes al ciclo abril--octubre de 2025 en Hidalgo, Puebla y Tlaxcala.

Este escenario se aborda como un problema transductivo. Es admisible estudiar los atributos de las 197 parcelas mediante transformaciones que no dependen del rendimiento; en cambio, ningún ajuste supervisado puede consultar las 59 etiquetas ocultas ni la etiqueta de una parcela retenida en validación. La transducción no es sinónimo de rellenar artificialmente los rendimientos faltantes ni de entrenar sobre 197 resultados conocidos. Su fundamento metodológico aparece en el aprendizaje semisupervisado y la regularización geométrica \parencite{Chapelle2006,Belkin2006,Zhou2004}.

La evidencia disponible comprende observaciones BASIC de Sentinel y Landsat, observaciones PRO de Planet, climatología y precipitación CHIRPS, productos topográficos, información de suelo y recursos hídricos. Las señales tienen distintas resoluciones, escalas, fechas y mecanismos de ausencia. Por ello, una cobertura extensa de imágenes no equivale a contar con una muestra supervisada grande; los 138 rendimientos siguen siendo la restricción estadística central \parencite{Funk2015CHIRPS,Poggio2021SoilGrids,INEGICEM,FAOWaPOR,SIAPDatos}.

La modelación requiere equilibrar flexibilidad e interpretabilidad. Los modelos locales reconocen heterogeneidad territorial, mientras las regresiones regularizadas disminuyen sensibilidad a covariables correlacionadas y tamaños muestrales limitados. Las comparaciones se realizaron mediante validación dejando fuera una parcela, con especial cuidado de que la selección de parámetros ocurra sin usar su rendimiento. No obstante, explorar muchas arquitecturas sobre las mismas 138 etiquetas puede introducir optimismo de selección \parencite{Stone1974,HoerlKennard1970}.

Se adoptó una arquitectura denominada Reference26, compuesta por un núcleo climático global, un experto geográfico local, una compuerta de composición geológica y una corrección residual. Posteriormente se evaluaron estrategias de imputación, mezclas con otros modelos y una corrección final mediante análogos de 26 atributos. Este documento explica tanto los mecanismos retenidos como las alternativas consideradas, detalla los resultados de las exclusiones auténticas y distingue las estimaciones de error observadas de la precisión todavía desconocida en las 59 parcelas objetivo.
